\documentclass{article}
\usepackage{/home/user1/code/tex/sty}
\usepackage{amsmath}
\usepackage{geometry}
\usepackage{hyperref}
\usepackage{graphicx}
\setlength{\parindent}{0pt}
\geometry{top=1in,bottom=1in,left=1in,right=1in}
\newcommand{\pic}[2]{\begin{center}\includegraphics[height=#1]{#2}\end{center}}
\begin{document}
Der Verbreiter der $phi^3$ Theorie enthält einen Logarithmus, der möglicherweise divergent ist. Seine Divergenz trifft in Abhängigkeit von die kinematische Bedingungen des Stoßes. Wenn er divergenb ist, müssen wir den Logarithmus in dem Gegenterm einschließen.
\end{document}
